\documentclass{situationist-preprint}

\title{Comprehensibility and Final Theories}
\msid{SCI-20-b}
\msfamily{SCI-20}
\msdiscipline{Philosophy of Science}
\mstype{Consequence}
\mskeywords{final theory; theory of everything; comprehensibility; best system; effective theories; understanding}

\begin{document}
\maketitle

\begin{abstract}
SCI-20-a argues that a final theory, a complete set of laws, probably exists, and that it may exceed human comprehension. We derive three consequences. The incomprehensibility of a final theory would not make it any less final, since finality is defined without reference to understanding. Human science could nonetheless possess comprehensible theories that are true approximations of the final theory in limited domains, so that the progress of physics does not depend on grasping it. And a final theory could be found without being recognized as final, so that the end of the search for laws need not be an event in the history of science.
\end{abstract}

\section{The premises}

By Definition~2.2 of SCI-20-a, a theory is final if it is true, stated in the vocabulary of the genuine magnitudes, finitely axiomatized, and entails every law of nature. It is a complete set of laws, not a description of a bottom level of objects; Lemma~3.1 of SCI-20-a shows that its existence does not require a fundamental level, which keeps it consistent with MET-10's thesis that every level is grounded in a lower one. Theorem~4.1 of SCI-20-a concludes that a final theory probably exists, and Proposition~5.1 that it is epistemically possible, and not improbable, that it is not humanly comprehensible in the sense of Definition~2.3. The reason is that the simplicity that makes a system best, on the account of laws assumed there \citep{Lewis1983,Lewis1994}, is measured in the vocabulary of nature, not in the cognitive costs of human beings.

\section{Consequences}

\begin{corollary}\label{c:indep}
Whether a final theory exists does not depend on whether it is humanly comprehensible.
\end{corollary}

\begin{proof}
None of the four clauses of Definition~2.2 of SCI-20-a mentions understanding, and the best-system criterion of Assumption~2.4 is fixed by the facts about the genuine magnitudes. The existence of a final theory is therefore settled by the world, and its comprehensibility by the relation between the world and human minds.
\end{proof}

\begin{corollary}\label{c:effective}
If a final theory exists but is not humanly comprehensible, there may still be humanly comprehensible theories that are true, as approximations, of limited domains, and whose laws are entailed by it as approximations valid in those domains.
\end{corollary}

\begin{proof}
A theory of a limited domain needs only the magnitudes and regularities relevant to that domain, and its comprehensibility depends on its own axioms, not on those of the theory that entails it. Nothing in the incomprehensibility of the final theory prevents its consequences for a domain from being summarized in a simple form, as the consequences of quantum electrodynamics for low-energy phenomena are summarized in classical electrodynamics.
\end{proof}

Corollary~\ref{c:effective} means that the progress of physics does not wait on the final theory. The theories in actual use are effective theories of limited domains, and the objection from effective field theory that SCI-20-a answers shows that this is how physics has advanced. The incomprehensibility of the final theory, if it obtains, would limit the unification of these theories in a human-graspable form, not their truth or their use.

\begin{corollary}\label{c:recognition}
A final theory could be found without being recognized as final.
\end{corollary}

\begin{proof}
Finality requires that the theory entail every law, and laws not yet discovered are among them. Evidence can show that a theory entails every law so far known and has no anomalies; it cannot show that no undiscovered law fails to follow from it. If the theory is in addition not humanly comprehensible, even its consequences for known laws may be checked only piecemeal.
\end{proof}

\citet{Weinberg1992} anticipated that a final theory would be recognized by its rigidity and simplicity. Corollary~\ref{c:recognition} implies that such marks are evidence of finality, not proof of it, and Proposition~5.1 of SCI-20-a implies that they may be invisible to us, since the relevant simplicity is measured in nature's vocabulary rather than ours.

\begin{objection}
A theory no human can understand is not a theory for us, and a final theory in that sense is of no scientific interest.
\end{objection}

\begin{reply}
Its interest does not depend on our grasping it. By Corollary~\ref{c:effective}, it makes our effective theories approximations of something true rather than of further approximations without end, and that is a claim about the structure of what science is investigating.
\end{reply}

\section{Conclusion}

A final theory may exist, be beyond human comprehension, and be approximated by comprehensible theories in every domain we investigate. Its existence is a fact about the laws; its comprehensibility, and even its recognition, are facts about us.

\begin{thebibliography}{99}\small
\bibitem[Lewis(1983)]{Lewis1983} Lewis, D. (1983). New work for a theory of universals. \emph{Australasian Journal of Philosophy}, 61(4), 343--377.
\bibitem[Lewis(1994)]{Lewis1994} Lewis, D. (1994). Humean supervenience debugged. \emph{Mind}, 103(412), 473--490.
\bibitem[Weinberg(1992)]{Weinberg1992} Weinberg, S. (1992). \emph{Dreams of a Final Theory}. Pantheon.
\end{thebibliography}

\end{document}
